\section{Giới thiệu}
Một yêu cầu như “thêm con mèo dưới gốc cây” xác định cả đối tượng cần thêm, vị trí tương đối và phần cảnh cần giữ. Đồ thị cảnh biểu diễn hai đối tượng bằng đỉnh và quan hệ có hướng bằng cạnh. Đồ án chọn nhánh ảnh tĩnh làm bài toán chính. Các mô hình thị giác và sinh ảnh được dùng ở chế độ
đóng băng; phần học là các mô hình trên đồ thị. Bảng~\ref{tab:scorecard} tổng hợp
kết quả trước khi đi vào thiết lập, đối chứng và phân tích lỗi.

\begin{table*}[!t]\centering\small
\caption{Kết quả chính của ba nhiệm vụ và thực nghiệm mở rộng. Các hàng khác dataset, split hoặc metric không dùng để xếp hạng chung.}\label{tab:scorecard}
\begin{tabularx}{\linewidth}{@{}p{10mm}p{28mm}p{29mm}X@{}}
\toprule NV & Phần & Dữ liệu / split & Kết quả chính\\\midrule
1 & Relation prediction & VG150 / test & GAT mR@50 \textbf{18,61\%}; hơn MLP \textbf{1,65 pp}.\\
1 & Loss ablation & VG150 / test & Adjusted, KG-off: mR@50 \textbf{24,59\%}, tăng \textbf{6,48 pp} so CE-off.\\
1 & KG fusion & VG150 / test & Pair-real mR@50 \textbf{25,24\%}; hơn shuffled 0,24 pp, chỉ 2/3 seed tăng.\\
2A & Phương pháp 2023+ & Custom GQA / test & PE-PLPR mR@50 \textbf{4,07\%}.\\
2A & AlignG ablation & Custom GQA / test & GC mR@50: \textbf{4,07 $\to$ 4,62\%}; bản adapted.\\
2B & Biến đổi đồ thị gốc & MagicBrush / test & GNN micro-F1 \textbf{7,34\%}; exact target graph \textbf{0\%}.\\
2B & Dataset thay thế & EMNLP 2020 / test & GNN micro-F1 \textbf{82,44\%}; macro-F1 \textbf{55,03\%}.\\
3 & Nhất quán đồ thị & VN100 / 200 edit & $\rho_G=\mathbf{0,1767}$; coverage \textbf{45,5\%}.\\
3 & Baseline ảnh & Tập giao $N=91$ & $-\mathrm{LPIPS}$ 0,1962; Graph 0,1767; CLIPScore 0,0561.\\
Video & Quan hệ video & PVSG / test & Conv kernel 5 đạt R@50 \textbf{8,63\%}; ST-GAT đạt mR@50 \textbf{2,38\%} trên track chuẩn.\\
\bottomrule\end{tabularx}
\end{table*}

\subsection{Quy ước đọc kết quả}
R và mR là phần trăm; pp là \emph{điểm phần trăm}, không phải phần trăm tăng
tương đối. Mean $\pm$ std là trung bình và độ lệch chuẩn mẫu qua ba seed,
không phải khoảng tin cậy. GC (graph constraint) giữ một predicate cho mỗi cặp
object có hướng; NG cho phép nhiều predicate. Bảng VG150 chính dùng GC
reference-v1 trên 26.446 ảnh test. Long-tail và zero-shot cũng đã bổ sung
kết quả test; custom GQA có cả tập validation và tập test của đồ án. Các ca định tính VG150 dùng validation.

\subsection{Ba kết luận cần kiểm chứng}
\begin{enumerate}
\item \textbf{Nhiệm vụ 1:} GAT có lợi so với MLP, nhưng điều chỉnh loss
tác động tới mR lớn hơn các cách thêm KG đã thử. Tăng mR không đồng nghĩa tăng R.
\item \textbf{Nhiệm vụ 2:} phương pháp mới tăng mR trên custom GQA; học graph operation
khả thi hơn khi supervision có cấu trúc. EMNLP không thay thế bằng chứng còn thiếu
của pipeline từ ảnh MagicBrush.
\item \textbf{Nhiệm vụ 3:} điểm đồ thị cung cấp chẩn đoán theo object/quan hệ nhưng
độ phủ thấp. Tương quan hiện tại chưa chứng minh ưu thế so với baseline ảnh.
\end{enumerate}
\srcnote{Nguồn số liệu: các report versioned trong \path{reports/level1}, \path{reports/level2}
và \path{docs/LEVEL3_RESULTS_V1.md}. Ánh xạ chi tiết nằm trong \path{SOURCE_MAP.md} đi kèm.}
